\documentclass[11pt]{article}
\usepackage[a4paper, margin=2.5cm]{geometry}
\usepackage[T1]{fontenc}
\usepackage{enumitem}
\usepackage{booktabs}
\usepackage{verbatim}
\usepackage{hyperref}

\title{Instruction for Reproducing Results\\[4pt]
\large COMP0197 Applied Deep Learning -- Team 09}
\author{}
\date{}

\begin{document}
\maketitle

\section{Environment}

All code was developed and tested under the \texttt{comp0197-pt} micromamba environment. The following three additional packages are required beyond the base environment:

\begin{table}[h]
\centering
\begin{tabular}{lll}
\toprule
Package & Version & Purpose \\
\midrule
\texttt{mne} & 1.6.0 & EEG data loading and filtering \\
\texttt{tqdm} & 4.66.1 & Progress bars during training/evaluation \\
\texttt{scikit-learn} & 1.3.2 & Metrics, calibration, and data splitting \\
\bottomrule
\end{tabular}
\end{table}

\noindent Install with:
\begin{verbatim}
pip install mne==1.6.0 tqdm==4.66.1 scikit-learn==1.3.2
\end{verbatim}

\section{Project Structure}

\begin{verbatim}
.
├── train.py                    # Main training entry point
├── test.py                     # Main evaluation entry point
├── saved_models/               # Trained model weights
│   └── best_model.pt
├── src/
│   ├── models/                 # CNN, CNN-LSTM, Deep Ensemble training
│   ├── preprocessing/          # Data pipeline, preprocessing, splitting
│   ├── uncertainty/            # MC Dropout, calibration, ensemble UQ
│   └── data/                   # Feature extraction and data loading
├── requirements.txt
└── instruction.pdf
\end{verbatim}

\section{Data Preparation}

\subsection{Step 1: Download CHB-MIT Dataset}

Download the raw EEG data (\textasciitilde42\,GB) from PhysioNet:

\begin{verbatim}
pip install awscli
aws s3 sync --no-sign-request \
    s3://physionet-open/chbmit/1.0.0/ \
    /path/to/data/raw/
\end{verbatim}

\subsection{Step 2: Inspect and Clean}

\begin{verbatim}
python src/preprocessing/data_pipeline.py inspect \
    --data-dir /path/to/data/raw
python src/preprocessing/data_pipeline.py clean \
    --data-dir /path/to/data/raw
\end{verbatim}

\subsection{Step 3: Preprocess}

\begin{verbatim}
python src/preprocessing/preprocess.py \
    --raw-dir /path/to/data/raw \
    --out-dir /path/to/data/processed \
    --stride-sec 30
\end{verbatim}

This produces 115{,}330 windows (30\,s each, 18 channels, 256\,Hz) stored as \texttt{.npy} files per patient.

\subsection{Step 4: Create Balanced Patient-level Split}

\begin{verbatim}
python src/preprocessing/make_balanced_split.py \
    --processed-dir /path/to/data/processed
\end{verbatim}

Output: \texttt{patient\_split\_balanced.json} (18 train / 3 val / 3 test patients).

\section{Reproducing Results}

\subsection{Step 5: Train the CNN+LSTM Model}

\begin{verbatim}
python train.py \
    --processed-dir /path/to/data/processed \
    --split-file /path/to/patient_split_balanced.json \
    --out-dir saved_models/ \
    --seq-len 6 --seq-stride 2 \
    --epochs 12 --batch-size 4 \
    --lr 1e-4 --weight-decay 1e-4 \
    --dropout 0.3 --patience 6
\end{verbatim}

The best checkpoint is saved as \texttt{saved\_models/best\_model.pt} based on validation AUPRC.

\subsection{Step 6: Evaluate on Test Set}

\begin{verbatim}
python test.py \
    --model-path saved_models/best_model.pt \
    --processed-dir /path/to/data/processed \
    --split-file /path/to/patient_split_balanced.json \
    --out-dir results/
\end{verbatim}

This script loads the saved model and produces:
\begin{itemize}[nosep]
    \item \texttt{test\_metrics.json} -- AUROC, AUPRC, F1, Brier, ECE (before and after calibration), MC Dropout error-detection scores.
    \item \texttt{test\_predictions.csv} -- Per-sample predictions.
    \item \texttt{roc\_pr\_curves.png} -- ROC and Precision-Recall curves.
    \item \texttt{reliability\_diagram.png} -- Calibration reliability diagram.
\end{itemize}

\subsection{Expected Test Results}

\begin{table}[h]
\centering
\begin{tabular}{lccccc}
\toprule
Pipeline & AUROC & AUPRC & F1 & Brier & ECE \\
\midrule
Tuned CNN+LSTM & 0.797 & 0.189 & 0.189 & 0.0602 & 0.0497 \\
+ Temperature Scaling & 0.797 & 0.189 & 0.189 & 0.0565 & 0.0133 \\
+ MC Dropout ($M$=20) & 0.785 & 0.182 & 0.203 & 0.0601 & -- \\
\bottomrule
\end{tabular}
\end{table}

\end{document}
